Por meio de condicionais e loops, é possível criar programas que executam 
seções específicas de código com base em determinadas condições definidas 
no próprio código.

\subsubsection{Condicionais}

Condicionais são estruturas de controle que direcionam o fluxo de execução 
do programa. Logo abaixo, vamos explorar de maneira simples como esses 
conceitos podem ser aplicados na criação de aplicações e por que eles são tão 
fundamentais:

\subsubsection{if}

A condicional \texttt{if} é a mais simples de toda a linguagem C e sua estrutura 
geralmente segue o seguinte formato:

\begin{excode}
    if (condição) {
        // Código a executar se condição for verdadeira...
    }
\end{excode}

A condição fornecida entre parênteses para o \texttt{if} deve ser uma 
condição que resulta em verdadeiro ou falso. Dependendo do resultado da 
condição, o bloco de código entre chaves será executado ou ignorado.

É possível utilizar valores constantes como \(0\) ou \(1\) para representar 
verdadeiro e falso, respectivamente. No entanto, é redundante utilizar uma expressão 
que sempre resulta em verdadeiro ou falso. Embora esse tipo de operação seja raramente 
utilizado, segue um exemplo para ilustrar o conceito:

\begin{excode}
    #include <stdio.h>

    int main(void)
    {
        if (0) { // 0 representa falso
            printf("Esse texto nunca será impresso\n");
        }
        if (1 > 2) { // 1 nunca será maior que 2
            printf("Esse texto nunca será impresso\n");
        }
        if (1) { // Qualquer número diferente de 0 representa verdadeiro
            printf("Esse texto sempre será impresso\n");
        }
        if (10 > 0) { // 10 sempre será maior que 0
            printf("Esse texto sempre será impresso\n");
        }
        
        return 0;
    }
\end{excode}

Neste exemplo, as verificações produzem resultados constantes, ou seja, o 
resultado não varia independentemente do número de vezes que o código é 
executado.

\subsubsection{else}

Outra palavra-chave importante de se entender é o \texttt{else}, que complementa o 
\texttt{if}. O bloco de código dentro do \texttt{else} é executado se a 
condição do \texttt{if} for falsa. Um exemplo prático ilustra sua utilização:

\begin{excode}
    #include <stdio.h>

    int main(void)
    {
        int number = 0;
        printf("Digite um número: ");
        scanf("\%d", &number);

        if (number > 0) {
            printf("Número é maior que 0\n");
        } else {
            printf("Número não é maior que 0\n");
        }
        
        return 0;
    }
\end{excode}

Neste código, se o número digitado for maior que \(0\), o primeiro bloco de 
código será executado; caso contrário, o bloco de código dentro do 
\texttt{else} será executado.

Além disso, é possível criar fluxos de execução mais complexos utilizando 
\texttt{if} dentro de \texttt{else} para definir condições adicionais. 
Apesar de ser complexo de entender apenas com explicações textuais, um 
exemplo de código esclarecerá o conceito:

\begin{excode}
    #include <stdio.h>

    int main(void)
    {
        int number = 0;
        printf("Digite um número: ");
        scanf("\%d", &number);

        if (number < 0) {
            // Executa bloco se número for menor que 0
            printf("Número é menor que 0\n");
        } else if (number > 0) {
            // Executa bloco se número for maior que 0
            printf("Número é maior que 0\n");
        } else {
            // Caso nenhuma das condições sejam
            // verdadeiras, o número só pode ser igual a 0
            printf("Número é exatamente igual a 0\n");
        }
        
        return 0;
    }
\end{excode}

Neste exemplo, além de verificar se o número é maior ou menor que \(0\), o 
programa também identifica quando o número é exatamente \(0\). Vale 
ressaltar que, ao utilizar \texttt{else if}, a expressão dentro dos 
parênteses só é avaliada se as condições anteriores forem falsas.

Em situações em que é necessário avaliar cada condição individualmente, 
vários \texttt{if} consecutivos são utilizados. Por exemplo:

\begin{excode}
    #include <stdio.h>

    int main(void)
    {
        int number = 120;

        if (number > 0) {
            printf("Número é maior que 0\n");
        }
        if (number > 10) {
            printf("Número é maior que 10\n");
        }
        if (number > 100) {
            printf("Número é maior que 100\n");
        }
        if (number > 1000) {
            printf("Número é maior que 1000\n");
        }

        return 0;
    }
\end{excode}

Ao executar o código acima, você deve ver algo como:

\begin{textbox}
    Número é maior que 0
    Número é maior que 10
    Número é maior que 100
\end{textbox}

Isso ocorre porque cada bloco \texttt{if} é independente dos outros; portanto, 
todos são avaliados individualmente. Se fosse utilizado um \texttt{else if}, como 
no exemplo a seguir:

\begin{excode}
    #include <stdio.h>

    int main(void)
    {
        int number = 120;

        if (number > 0) {
            printf("Número é maior que 0\n");
        } else if (number > 10) {
            printf("Número é maior que 10\n");
        } else if (number > 100) {
            printf("Número é maior que 100\n");
        } else if (number > 1000) {
            printf("Número é maior que 1000\n");
        }

        return 0;
    }
\end{excode}

O resultado seria apenas:

\begin{textbox}
    Número é maior que 0
\end{textbox}

Isso ocorre porque a primeira condição já é avaliada como verdadeira, e as 
condições subsequentes dependem da falsidade das condições anteriores, o que 
não ocorre neste caso.

\subsubsection{switch}

Semelhante ao \texttt{if} concatenando uma sequência de \texttt{else} e \texttt{if}. 
O \texttt{switch} existe para executar um bloco de código específico em um cenário de 
múltiplas possibilidades para uma expressão. Teorias à parte, a sua estrutura parece 
algo como:

\begin{excode}
    switch (expressão) {
        case valor1:  
            // bloco de código 1
            break;
        case valor2:
            // bloco de código 2
            break;
        default:
            // código a rodar caso nenhuma das alternativas acima sejam válidas
            // Geralmente 'default' é a última seção por padrão
    }
\end{excode}

Neste caso, o código deve avaliar a expressão como algum valor inteiro ou até mesmo 
um valor de caractere, esse valor então será comparado com cada um dos valores 
especificados após os \texttt{case} que são, de forma simplificada, uma seção de 
código que devem rodar em caso da expressão ser igual ao valor passado, e caso o 
valor da expressão não seja igual a nenhum dos valores especificados a seção 
definida por \texttt{default} executará por padrão. Após cair em uma seção em 
específico, o código executará até que atinja um \texttt{break}, que é usado 
justamente para delimitar todo o código que deve ser executado.

Dito isso tudo, vamos para alguns exemplos práticos que podem ajudar a entender melhor como 
isso tudo funciona:

\begin{excode}
    #include <stdio.h>

    int main(void)
    {
        int dia = 5; // Mude à vontade

        // Associa o valor do dia com os dias da semana
        switch (dia) {
            case 1:
                printf("Segunda\n");
                break;
            case 2:
                printf("Terça\n");
                break;
            case 3:
                printf("Quarta\n");
                break;
            case 4:
                printf("Quinta\n");
                break;
            case 5:
                printf("Sexta\n");
                break;
            case 6:
                printf("Sábado\n");
                break;
            case 7:
                printf("Domingo\n");
                break;
            default:
                printf("Dia inválido\n");
        }

        return 0;
    }
\end{excode}

No caso acima temos um exemplo bem simples, de como usar o switch. Utilizamos o valor 
\texttt{dia} como expressão e caso o seu valor seja igual a algum dos especificados ele 
deve executar a sua respectiva seção. Nesse caso em específico, ao executar o código teremos 
algo como:

\begin{textbox}
    Sexta
\end{textbox}

É notável também que a seção \texttt{default} é usada para especificar qual comportamento 
deve ocorrer caso o valor \texttt{dia} não corresponda a nenhuma das outras seções, exibindo 
nesse caso que não foi possível identificar o número.

\begin{postscript}
    Outro detalhe importante, é que é possível atribuir vários casos a uma mesma seção de 
    código, sendo possível fazer algo como por exemplo:

    \begin{excode}
    #include <stdio.h>

    int main(void)
    {
        int dia = 5; // Mude à vontade

        switch (dia) {
            case 1: case 2: case 3: case 4: case 5:
                printf("Dia útil\n");
                break;
            case 6:
            case 7:
                printf("Fim de semana\n");
                break;
            default:
                printf("Dia inválido\n");
        }

        return 0;
    }
    \end{excode}

    Aqui, combinamos vários casos para verificar se o valor \texttt{dia} corresponde a um 
    dia útil ou a um final de semana. Sendo possível dizer que o \texttt{switch} é semelhante 
    a uma sequência de \texttt{if} e \texttt{else if}, usado quando há várias possibilidades 
    para uma expressão.
\end{postscript}

\subsubsection{while}

Essa palavra-chave, além de ser a estrutura de loop mais simples da linguagem, é 
responsável por definir a estrutura de repetição de uma seção de código enquanto uma dada
expressão for verdadeira.

\begin{excode}
    while (expressão) {
        // Bloco de código a ser executado
    }
\end{excode}

Vamos ilustrar um simples caso de uso para a estrutura while:

\begin{excode}
    #include <stdio.h>

    int main(void)
    {
        int x = 0;

        while (x != 10) {
            printf("x: %d\n", x);
            ++x; // Incrementa o valor de x em 1
        }

        return 0;
    }
\end{excode}

Neste exemplo, criamos uma estrutura \texttt{while} que verifica se \(x\) é 
diferente de 10. Enquanto o valor for diferente de 10, o código entre chaves é 
repetido. Quando executado, este código produzirá a seguinte saída:

\begin{postscript}
    É importante mencionar que o loop \texttt{while} está verificando apenas se 
    o valor é diferente de 10, o que poderia ocasionar erros caso fosse 
    possível o valor de \texttt{x} pular de 9 para 11, caso ele fosse 
    incrementado de 2 em 2 por exemplo.

    No caso acima isso não é possível, mas deve-se ficar atento a casos 
    específicos como esse em aplicações reais que usem algum loop do tipo, 
    em caso de dúvida vale sempre a pena usar comparações como:

    \begin{excode}
    while (x < 10) {
        ...
    }
    \end{excode}

    Assim, você garante a quebra do loop quando o valor de \texttt{x} for maior 
    ou igual a 10.
\end{postscript}

\begin{textbox}
    x: 0
    x: 1
    x: 2
    x: 3
    x: 4
    x: 5
    x: 6
    x: 7
    x: 8
    x: 9
\end{textbox}

Essa funcionalidade pode ser usada para criar algoritmos mais complexos, 
repetindo um mesmo código várias vezes sem a necessidade de escrever mais 
linhas de código.

Um código equivalente, porém menos prático de se escrever e manter, seria:

\begin{excode}
    #include <stdio.h>

    int main(void)
    {
        printf("x: 0\n");
        printf("x: 1\n");
        printf("x: 2\n");
        printf("x: 3\n");
        printf("x: 4\n");
        printf("x: 5\n");
        printf("x: 6\n");
        printf("x: 7\n");
        printf("x: 8\n");
        printf("x: 9\n");

        return 0;
    }
\end{excode}

Imagine repetir essa estrutura para números maiores, como até 10000. Seria um 
processo extremamente tedioso e ineficiente ter que escrever tudo isso 
manualmente, daí a utilidade das estruturas como \texttt{while}.

\subsubsection{for}

Outra estrutura de loop semelhante, mas um pouco mais complexa, é o 
\texttt{for} (conhecido também como for loop). Esta estrutura segue o seguinte padrão:

\begin{excode}
    for (inicialização; condição; expressão) {
        // Código a ser executado
    }
\end{excode}

Dentro de uma única estrutura, há três seções com características únicas:

\begin{itemize}
    \item \textbf{Inicialização}: uma expressão que é executada uma única vez 
        antes do loop começar.
    \item \textbf{Condição (de parada)}: verificada antes de cada execução do 
        loop.
    \item \textbf{Expressão}: geralmente é um incremento, mas pode ser qualquer 
        tipo de expressão que precise ser executada ao fim do loop.
\end{itemize}

Para ilustrar, vamos usar essa estrutura para replicar o exemplo do \texttt{while}:

\begin{excode}
    #include <stdio.h>

    int main(void)
    {
        for (int x = 0; x != 10; ++x) {
            printf("x: \%d\n", x);
        }

        return 0;
    }
\end{excode}

Para entender melhor o código acima, vamos explicar cada parâmetro 
individualmente:

\begin{itemize}
    \item \texttt{int x = 0}: Inicializa a variável \(x\) com o valor 0.
    \item \texttt{x != 10}: Expressão avaliada no início de cada loop.
    \item \texttt{++x}: Incrementa o valor de \(x\) ao final de cada loop.
\end{itemize}


Ao executar o código, o resultado será idêntico ao exemplo utilizando 
\texttt{while}.

É importante notar que, em um \texttt{for} loop, é possível omitir qualquer 
um desses parâmetros, embora seja contra-intuitivo. Por exemplo, podemos criar 
um loop infinito da seguinte forma:

\begin{excode}
    #include <stdio.h>

    int main(void)
    {
        for (;;) {
            // Esse loop deve continuar sendo executado indefinidamente 
        }

        while (1) {
            // O while requer uma expressão qualquer, desde que a mesma 
            // diferente de zero para funcionar como um loop infinito
        }

        return 0;
    }
\end{excode}

Sabendo que é possível omitir parâmetros de um \texttt{for}, é possível 
também, imitar o funcionamento de um \texttt{while} loop usando um 
\texttt{for} loop da seguinte forma:

\begin{excode}
    #include <stdio.h>

    int main() {
        int i = 0;
        for (; i < 5; ) { // inicialização vazia, condição igual ao while, expressão vazia
            printf("%d\n", i);
            i++; // incremento feito no corpo
        }
        return 0;
    }
\end{excode}

Fazendo um \texttt{for} com a inicialização e a expressão vazia, podemos 
imitar perfeitamente o funcionamento de um \texttt{while}.

\subsubsection{do while}

Também temos o loop menos utilizado da linguagem C, chamado 
\texttt{do while}. A diferença fundamental deste loop é que ele executa o 
bloco de código antes de verificar a expressão. A estrutura deste loop é a 
seguinte:

\begin{excode}
    do {
        // Código a ser executado
    } while (expressão);
\end{excode}

Vamos demonstrar o funcionamento desse loop com um exemplo simples:

\begin{excode}
    #include <stdio.h>

    int main(void)
    {
        do {
            printf("Esse texto será impresso\n");
        } while (0); // Verificação é feita após a execução

        while (0) { // Verificação é feita antes da execução
            printf("Esse texto não será impresso\n");
        }

        return 0;
    }
\end{excode}

Caso executemos o código acima, veremos o seguinte resultado:

\begin{textbox}
    Esse texto será impresso
\end{textbox}

Nota-se que o texto do loop \texttt{while}, ao contrário do texto do loop 
\texttt{do while}, não foi impresso. Isso ocorre porque o \texttt{do while} 
primeiro executa o código dentro do bloco e, em seguida, verifica a expressão. 
Como ambas as expressões no código acima são falsas, o \texttt{do while} 
executará o código uma vez e sairá do loop, enquanto o \texttt{while} 
nunca executará seu bloco de código.

Novamente, é possível imitar o funcionamento de um \texttt{for} 
loop usando o \texttt{do while}:

\begin{excode}
    do
    {
        //  Código a ser executado
        printf( "%d ", i);   
        // Expressão
        i++;
    }  while (i <= 10); // Checagem de condição 
\end{excode}

O código acima deve executar de maneira quase equivalente ao \texttt{for} loop, com a 
única exceção de que a primeira iteração será executada antes da checagem da condição 
ser feita.

\subsubsection{Exercícios}

\textbf{É de extrema importância realizar os exercícios aqui elaborados 
para fixar o conteúdo.}

\begin{enumerate}
    \item {Escreva um programa que verifica se um número inteiro e informe se ele é 
        positivo, negativo ou zero. Utilize as estruturas \texttt{if}, 
        \texttt{else if} e \texttt{else}.}
    \item {Modifique o valor do dia nos códigos de exemplo do \texttt{switch} e 
        observe o que acontece}
    \item {Crie um programa que imprima todos os números pares de 1 até 100 
        utilizando um loop \texttt{while}. Em seguida, refaça o mesmo programa 
        usando um loop \texttt{for}.}
    \item {Escreva um código que leia um caractere. Se o caractere for uma vogal 
        (a, e, i, o, u, maiúscula ou minúscula), exiba uma mensagem informando 
        que é vogal. Caso contrário, informe que não é vogal. Utilize \texttt{if} 
        e \texttt{else}.}
    \item {Crie um programa que exiba a tabuada de um número (de 1 a 10). Utilize 
        um loop \texttt{for} para gerar e mostrar cada linha da tabuada no formato: 
        \texttt{5 x 1 = 5}.}
\end{enumerate}
